% ===================================================================
%  அட்டவணை எண் 12 — நிலையம். இரயில் பாதைகள். கப்பல்கள்.
%
%  Traduction, faite en 2026, en tamoul d'aujourd'hui, sur l'ido de
%  texto/io. Regles de la colonne : voir 10-attavanai-01.tex. Recit a
%  la premiere personne, sans scenes. Ecrit avec outils/ordo.py sous
%  les yeux, bloc par bloc.
%
%  LE CHEMIN DE FER NE COUTE RIEN AU TAMOUL, ET LA DATE LE DIT. La
%  premiere ligne indienne part de Madras en 1856 ; la langue a donc
%  soixante-dix ans de vocabulaire ferroviaire quand ce livret
%  parait, et il est a elle : தண்டவாளம் le rail, நடைமேடை le quai,
%  பெட்டி la voiture, கொதிகலன் la chaudiere, சக்கரம் la roue, விசில்
%  le sifflet, சுரங்கப் பாதை le tunnel, சரக்கு இரயில் le train de
%  marchandises, முள் மாற்றுபவர் l'aiguilleur, சமிக்ஞை le signal.
%  C'est le troisieme tableau d'affilee — apres le port et la ville —
%  ou la colonne n'emprunte presque rien, et le tableau 9 reste seul
%  a son extreme. La colonne telougoue avait fait exactement le meme
%  parcours, et pour la meme raison : le rail est entre en Andhra en
%  1893.
%
%  MAIS DEUX MOTS DE FRANCE RESTENT DES MOTS DE FRANCE. Le bloc 19
%  nomme « Deutsch-Avricourt », la gare-frontiere de 1926, et la note
%  (*) precise que la description suit la frontiere d'avant-guerre.
%  Ce n'est pas une chose a traduire : c'est un lieu et une date. La
%  colonne translittere le nom et rend la note telle quelle.
%
%  LE FAC-SIMILE COUPE LE BLOC 14 EN DEUX PAR SA NOTE, et la
%  traduction fait de meme : t12-14-1, puis t12-noto-1, puis la suite
%  de t12-14-1. Les renvois y courent 64, 65, 66 — note — 63, 67, 68,
%  69. Le (63) tombe donc APRES le (66), ce qui a l'air d'une faute
%  et n'en est pas une : c'est la voie d'a cote, nommee quand le train
%  la depasse. ordo.py recolle les deux moities pour donner la suite
%  vraie, et c'est exactement ce pour quoi il a ete ecrit.
%
%  ET LA « NOTO » DE L'OUVERTURE APPARTIENT A L'APPARAT, PAS AU
%  PREMIER ALINEA. Je l'avais mise dans t12-01-1 ; kolonoj.py a
%  signale t12-apar-1 tombe a 36 % de la longueur mediane. Le
%  fac-simile la compose sous le titre et avant le « 1. --- », donc
%  dans le bloc d'apparat, et la colonne coreenne l'y avait deja
%  mise. Un bloc trop court est presque toujours un bloc dont le
%  contenu est parti chez le voisin.
%
%  TRENTE PAIRES, DIX-SEPT RETOURNEMENTS. Le tableau attache
%  beaucoup — l'employe COMMIS a la bascule, le chauffeur SUR son
%  tender, la chaudiere DE la locomotive, le pont QUI franchit le
%  fleuve — mais la moitie du texte enumere les guichets d'une gare,
%  et l'enumeration ne coute rien.
% ===================================================================

%%K t12-apar-1 apar
\VUsaut{4.0mm}
\VUtitre{15.0pt}{60}{அட்டவணை எண் 12}
\VUsaut{2.4mm}
\VUpk{11.4pt}{40}{நிலையம். --- இரயில் பாதைகள். --- கப்பல்கள்.}
\VUsaut{3.4mm}
\textsc{குறிப்பு.} --- \textit{ஒவ்வொரு நாட்டிலும் உபயோகிக்கும் நேர}
\textit{அட்டவணைகள் வெவ்வேறாக இருக்கின்றன ; அதனால் எடுத்துக்காட்டாக}
\VUgras{பிரெஞ்சு அட்டவணையின்} \textit{நேரங்களை உபயோகிக்கிறோம்.}

%%K t12-01-1 p
1. --- நான் இப்போதுதான் ஜெர்மனி வழியாகப் பயணம் செய்தேன். புறப்படுவதற்கு
முன், இரயில்களின் புறப்படும் நேரத்தை அறிய ஒரு \VUgras{கால அட்டவணைப்
புத்தகத்தை} \textsuperscript{(15)} வாங்கினேன். மிகவும் வசதியான இரயில்
பாரிஸிலிருந்து மாலை ஒன்பது மணிக்குப் புறப்பட்டு, ஸ்ட்ராஸ்பூர்குக்கு
ஒரு மணி பதின்மூன்று நிமிடத்திற்கு வரும் என்று உறுதி செய்த பிறகு, என்
பயணத்தைத் தயார் செய்து, மறு நாள் நிலையத்திற்குக் கொண்டு போகச்
சொன்னேன்.

%%K t12-02-1 p
2. --- பெரிய புறப்படும் கூடத்திற்குள் நான் நுழைந்த போது, ஒரு நாட்டுப்புற
\VUgras{குருவுக்குப்} \textsuperscript{(2)} பின்னால் — அவர் கையில் தன்
\VUgras{பயணப் பையைச்} \textsuperscript{(3)} சுமந்து வந்தார் —, ஏராளமான
\VUgras{பயணிகள்} \textsuperscript{(1)} ஏற்கெனவே வந்திருந்தார்கள்.

%%K t12-02-2 p
ஒரு \VUgras{தந்தி} \textsuperscript{(5)} அனுப்ப விரும்பியதால், ஒரு
\VUgras{சோதனையாளரிடம்} \textsuperscript{(6)} \VUgras{தந்தி அலுவலகத்தைக்}
\textsuperscript{(4)} காட்டச் சொன்னேன்.

%%K t12-03-1 p
3. --- \VUgras{காத்திருப்பு அறைக்குள்} \textsuperscript{(7)} போவதற்கு
முன், திரும்பி வரும் \VUgras{பயணச் சீட்டை} \textsuperscript{(9)} வாங்கப்
போனேன்.

%%K t12-04-1 p
4. --- \VUgras{சீட்டுச் சாளரத்தில்} \textsuperscript{(8)} நான் நெடு நேரம்
காத்திருக்கவில்லை ; ஏனெனில் அங்கே சிலரே இருந்தார்கள். விடுமுறையில்
புறப்படும் ஒரு \VUgras{வீரன்} \textsuperscript{(10)} தன் முறையை எனக்குத்
தர இணங்கினான் ; அதனால் \VUgras{நிலைய அதிகாரியிடம்} \textsuperscript{(13)}
சில செய்திகளைக் கேட்க எனக்குப் போதுமான நேரம் கிடைத்தது ; அவர்
கண்காணிப்பின் \VUgras{ஆணையருடனும்} \textsuperscript{(12)}, பணியில்
இருந்த \VUgras{படைக் காவலருடனும்} \textsuperscript{(11)} பேசிக்
கொண்டிருந்தார்.

%%K t12-tit-1 sub
\VUpk{10.2pt}{40}{பொதிப் பதிவு.}
\VUsaut{3.0mm}

%%K t12-05-1 p
5. --- \VUgras{புத்தகக் கடையில்} \textsuperscript{(14)} செய்தித்தாள்
ஒன்றை வாங்கிய பிறகு, என் \VUgras{பயணப் பொதிகளை}
\textsuperscript{(17)} எடை போடச் சொன்னேன் ; அவற்றை ஒரு \VUgras{சுமை
தூக்குபவர்} \textsuperscript{(16)} \VUgras{பொதி வண்டியில்}
\textsuperscript{(19)} இப்போதுதான் கொண்டு வந்தார். \VUgras{எழுத்தர்} \textsuperscript{(22)} — அவர் \VUgras{எடை போடும்
தராசுக்குப்} \textsuperscript{(21)} பொறுப்பானவர் — என் பெட்டகம் முப்பத்தைந்து கிலோ எடை என்று
எனக்குத் தெரிவித்தார் ; அது ஐந்து கிலோ \VUgras{கூடுதல் எடை} ;
அதற்காக \VUgras{பொதிச் சாளரத்தில்} \textsuperscript{(23)} கூடுதலாகச்
செலுத்த வேண்டும்.

%%K t12-06-1 p
6. --- நான் மிகவும் முன்னதாக வந்திருந்ததால், நிலையத்தில் பயணிகளின்
நடமாட்டத்தைக் கவனிக்க எனக்கு ஓய்வு நேரம் இருந்தது. சிலர் தங்கள்
சிறு பொதிகளை \VUgras{பொதி வைப்பகத்தில்} \textsuperscript{(25)}
வைத்தார்கள் அல்லது திரும்ப எடுத்தார்கள் ; வேறு சிலர் \VUgras{கால
அட்டவணைகளைப்} \textsuperscript{(26)} பார்த்தார்கள். வந்திறங்கிய
பயணிகள் \VUgras{வெளியேறும் வழிக்கு} \textsuperscript{(32)}
விரைந்தார்கள் — கையில் தங்கள் \VUgras{கைப் பெட்டியுடன்}
\textsuperscript{(24)}.
சிலர் \VUgras{பொதி வழங்கும் இடத்தில்} \textsuperscript{(27)}
காத்திருந்து, தங்கள் \VUgras{சீட்டைக்} \textsuperscript{(29)}
காட்டினார்கள் ; தங்கள் பெட்டகங்களையும் \VUgras{பெட்டிகளையும்}
\textsuperscript{(28)} \VUgras{கூடைகளையும்} \textsuperscript{(30)}
திரும்பப் பெறுவதற்காக.

%%K t12-07-1 p
7. --- \VUgras{கலால் வரி அலுவலர்கள்} \textsuperscript{(33)} பொதிகளைக்
கவனமாகச் சோதித்தார்கள் ; அறிவிக்க வேண்டியது ஏதேனும் இருக்கிறதா
என்று உறுதி செய்வதற்காக.

%%K t12-08-1 p
8. --- சில பயணிகள் \VUgras{உணவுக் கூடத்திற்குப்} \textsuperscript{(34)}
போனார்கள் ; அங்கே ஒரு \VUgras{பணியாளன்} \textsuperscript{(35)}
அவர்களுக்குப் பரிமாற ஆர்வம் காட்டினான். அவர்கள் அவசரப்பட வேண்டும் ;
ஏனெனில் விரைவு \VUgras{இரயில்} \textsuperscript{(36)} நிலையத்தில் நெடு
நேரம் நிற்காது.

%%K t12-09-1 p
9. --- பிறகு புறப்படும் நடைமேடைக்குப் போய், நேரடி \VUgras{இரயில்
பெட்டியைத்} \textsuperscript{(37)} தேடினேன் ; பயணத்தின் நடுவில் பெட்டி
மாற வேண்டாம் என்பதற்காக. நடைபாதை உள்ள ஒரு பெட்டியில் ஏறினேன் ; அதில்
புகைப்பவர்களுக்கு ஒரு \VUgras{பயண அறையும்} \textsuperscript{(38)},
« தனியாக வரும் பெண்களுக்கு » ஒதுக்கிய அறையும் இருந்தன.

%%K t12-tit-2 sub
\VUpk{10.2pt}{40}{இரவில் புறப்பாடு.}
\VUsaut{3.0mm}

%%K t12-10-1 p
10. --- \VUgras{கால் படியில்} \textsuperscript{(40)} ஏறி,
\VUgras{கதவை} \textsuperscript{(39)} அடைத்து, \VUgras{திரையைச்}
\textsuperscript{(41)} சுருட்டி, என் சிறு பொதிகளை \VUgras{வலை
அடுக்கின்} \textsuperscript{(43)} மேல் வைத்த பிறகு,
\VUgras{இருக்கையில்} \textsuperscript{(42)} அமர்ந்தேன்.
\VUgras{விளக்கு ஏற்றுபவர்} \textsuperscript{(44)} விளக்குகளை ஏற்றுவதைப்
பார்த்து, இரவைப் பெட்டியில் கழிக்க வேண்டும் என்று எண்ணி,
\VUgras{தலையணைகளை} \textsuperscript{(45)} வாடகைக்குத் தரும்
அலுவலரை அழைத்து ஒன்று கேட்டேன்.

%%K t12-11-1 p
11. --- இரயில் நடத்துநர் சீட்டுகளைச் சரி பார்க்க வந்தார். சரியான
நேரம் ஆகியும் நாம் ஏன் இன்னும் புறப்படவில்லை என்று அவரிடம்
கேட்டேன். \VUgras{இரயில் மேற்பார்வையாளர்} \textsuperscript{(46)}
\VUgras{சரக்குப் பெட்டியில்} \textsuperscript{(47)} மிதிவண்டிகளை
வைப்பதில் ஈடுபட்டிருக்கிறார் என்று பதில் சொன்னார். மேலும்
கால் சூடேற்றிகள் \textsuperscript{(48)} எல்லாம் இன்னும்
மாற்றப்படவில்லை.

%%K t12-12-1 p
12. --- ஆனால் நாம் விரைவில் புறப்படவிருந்தோம் ; ஏனெனில்
\VUgras{நெருப்பு ஏற்றுபவர்} \textsuperscript{(50)} — அவர்
தன் \VUgras{நிலக்கரிப் பெட்டியின்} \textsuperscript{(49)} மேல்
இருந்தார் — \VUgras{இரயில்
இயந்திரத்தின்} \textsuperscript{(54)} நெருப்பைத் தூண்ட நிலக்கரி
எடுத்தார் ; \VUgras{இயந்திர ஓட்டுநர்} \textsuperscript{(51)}
\VUgras{நிலைய உதவி அதிகாரியின்} \textsuperscript{(52)} சமிக்ஞைக்கு
மட்டுமே காத்திருந்தார் ; அந்த அதிகாரி \VUgras{நடைமேடையின்}
\textsuperscript{(53)} மேல் இருந்தார்.

%%K t12-13-1 p
13. --- இரயில் இயந்திரத்தின் \VUgras{கொதிகலன்} \textsuperscript{(56)}
மந்தமாக முழங்கியது ; \VUgras{புகைபோக்கி} \textsuperscript{(55)}
\VUgras{நீராவியுடன்} \textsuperscript{(60)} கலந்த புகையை
வெள்ளமாகக் கக்கியது. கூர்மையான \VUgras{விசில்} \textsuperscript{(59)}
ஒலித்தது ; \VUgras{உந்து தண்டுகள்} \textsuperscript{(58)} அசையத்
தொடங்கி, அந்தப் பாரமான இயந்திரத்தின் \VUgras{சக்கரங்களைத்}
\textsuperscript{(57)} தள்ளின.

%%K t12-tit-3 sub
\VUsaut{3.0mm}
\VUpk{10.2pt}{40}{புறப்பாடு!}
\VUsaut{3.0mm}

%%K t12-14-1 p
14. --- \VUgras{தண்டவாளங்களின்} \textsuperscript{(64)} மேல் ஓடும்
இரயிலின் வேகம் கூடியது ; \VUgras{நடைமேடைகள்} \textsuperscript{(65)}
மறைந்தன ; \VUgras{கழிப்பறைகளைக்} \textsuperscript{(66)}
கடந்தோம் ; பெட்டி வரிசை ஒன்றையும் கடந்தோம்.

%%K t12-noto-1 noto
\VUnotes{143.2mm}{%
\textit{(*) குறிப்பு.} --- \textit{இந்தப் பயண விளக்கம் போருக்கு}
\textit{முந்திய எல்லையின்படி எழுதப்பட்டது என்பது சொல்லாமலே}
\textit{விளங்கும்.}%
}

%%K t12-14-1 p suite
அந்தப் பெட்டி வரிசை மறு \VUgras{பாதையில்} \textsuperscript{(63)}
நிறுத்தப்பட்டிருந்தது : \VUgras{படுக்கைப் பெட்டிகள்}
\textsuperscript{(67)}, \VUgras{உணவுப் பெட்டி} \textsuperscript{(68)},
\VUgras{அஞ்சல் பெட்டி} \textsuperscript{(69)}.

%%K t12-15-1 p
15. --- பிறகு \VUgras{சரக்கு நிலையம்} \textsuperscript{(71)} நம்
கண்களுக்கு முன்னால் கடந்தது. கொட்டகைகளின் கீழ் அலுவலர்கள்
நிதானமான இரயிலில் அனுப்பும் \VUgras{சரக்குகளை} \textsuperscript{(72)}
ஏற்றினார்கள். \VUgras{பணிக் குழுவினர்} \textsuperscript{(76)}
\VUgras{நீர்த் தேக்கத்திற்கு} \textsuperscript{(73)} அருகில்
\VUgras{கால்நடைப் பெட்டிகளையும்} \textsuperscript{(75)}
\VUgras{தட்டைப் பெட்டிகளையும்} \textsuperscript{(79)} நகர்த்தி
\VUgras{சரக்கு இரயிலை} \textsuperscript{(80)} அமைத்தார்கள்.

%%K t12-16-1 p
16. --- கடைசி \VUgras{பாதை பிரியும் இடத்தைக்} \textsuperscript{(78)}
கடந்தோம் ; அதற்கு அருகில் முள் மாற்றுபவர் நின்றார்.
\VUgras{சமிக்ஞைத் தட்டைக்} \textsuperscript{(86)} கடந்தோம் ;
நிலையம் தொலைவில் மறைந்தது.

%%K t12-17-1 p
17. --- விரைவில் ஓர் \VUgras{இரும்புப் பாலத்தைக்} \textsuperscript{(74)}
கடந்தோம் ; அது \VUgras{பேராற்றைத்} \textsuperscript{(81)} தாண்டுகிறது.
அந்தப் பாலத்தைக் கடந்த பிறகு ஆற்றின் ஓரமாகச் சென்றோம் ; அதன் மேல்
வலிமையான \VUgras{இழுவைப் படகுகள்} \textsuperscript{(82)} பாரமான
\VUgras{சரக்குப் படகுகளை} \textsuperscript{(83)} இழுத்தன. மேலும் ஒரு
\VUgras{பயணக் கப்பலைப்} \textsuperscript{(84)} பார்த்தோம் ; அது ஒரு
\VUgras{மிதவைத் தளத்தை} \textsuperscript{(85)} அணுகிக் கொண்டிருந்தது.

%%K t12-18-1 p
18. --- பல நிலையங்களை மிக வேகமாகக் கடந்த பிறகு, சில \VUgras{சுரங்கப்
பாதைகளைக்} \textsuperscript{(87)} கடந்து, தடுப்புக் காவலர்களின்
எண்ணற்ற \VUgras{சிறு வீடுகளை} \textsuperscript{(88)} பின்னால் விட்டோம்.
இரயிலின் ஆட்டத்தால் தாலாட்டப்பட்டு ஆழ்ந்து தூங்கத் தொடங்கினேன்.

%%K t12-tit-4 sub
\VUsaut{3.0mm}
\VUpk{10.2pt}{40}{டாய்ச்-அவ்ரிகூர் நிலையம்.}
\VUsaut{3.0mm}

%%K t12-19-1 p
19. --- « டாய்ச்-அவ்ரிகூர்! \textsuperscript{(*)} எல்லாரும்
இறங்குங்கள்! » என்று கத்திய அலுவலரின் குரலால் திடீரென்று
எழுப்பப்பட்டேன். சுங்கச் சோதனையும் எல்லையின் அலுப்பூட்டும்
முறைமைகள் அனைத்தும் நடந்தன. என் சட்டைப் பை கடிகாரத்தை நிலையத்தின்
\VUgras{பெருங் கடிகாரத்திற்கு} \textsuperscript{(89)} ஒப்ப வைக்க
வேண்டியிருந்தது ; அது ஐம்பத்தைந்து நிமிடம் முன்னதாக இருந்தது.
அரிதாகவே விழித்திருந்த நான் \VUgras{பெரிய முள்ளையும்}
\textsuperscript{(90)} \VUgras{சிறிய முள்ளையும்} \textsuperscript{(91)}
சிரமப்பட்டே பிரித்து அறிந்தேன் ; காலைப் பனி மூட்டத்தால்
\VUgras{எண் தட்டே} \textsuperscript{(92)} முற்றிலும் தெளிவற்று
இருந்தது.

%%K t12-20-1 p
20. --- சுங்க அலுவலர்கள் சோதனை செய்யும் போது, நான் கொட்டாவி
விட்டுக் கொண்டே நிலையச் சுவர்களை அலங்கரிக்கும் \VUgras{விளம்பரத்
தாள்களைப்} \textsuperscript{(93)} படித்தேன். நேரத்தைக் கழிக்கக் காலை
உணவு சாப்பிட்டேன் ; ஸ்ட்ராஸ்பூர்குக்கு இன்னும் எவ்வளவு தூரம் என்று
பார்க்க ஜெர்மனிய \VUgras{வலையமைப்பின்} \textsuperscript{(94)}
வரைபடத்தைப் பார்த்தேன் ; என் பயணத்தின் இலக்கை விரைவில் அடையும்
நம்பிக்கையால் வலுப்பெற்று என் பெட்டிக்குள் மீண்டும் நுழைந்தேன்.

\VUsaut{6.0mm}
\VUfilet{22mm}
